\section{Eleven missing four-terminal sectors for seven ports}
\label{sec:path-seven-grid-count}

The weighted $K_{3,3}$ argument has a stronger counting form for seven ports.
Only four-terminal sectors are needed for the count below.

\begin{lemma}[Seven tokens in a three-by-three grid]
\label{lem:path-seven-grid-eleven}
Place seven distinct tokens in the cells of a $3\times3$ grid, allowing repeated cells.
At least eleven four-token subsets have all their cells contained in the union of one row and one column.
The bound is attained by placing $3,2,2$ tokens in three cells forming a perfect matching.
\end{lemma}

\begin{proof}
Call a four-token subset bad if it lies in such a row-column union, and good otherwise.
There are $\binom74=35$ subsets.
We prove that at most twenty-four are good.

If a row contains $r\ge4$ tokens, every subset with at least three tokens from that row is bad.
This gives
\[
 \binom r4+\binom r3(7-r)\ge13
 \qquad(4\le r\le7).
\]
We may therefore assume that the row counts are $(3,3,1)$ or $(3,2,2)$.

Here is a formula for the number $G$ of good subsets.
Write $w_{ic}$ for the token count in row $i$, column $c$, and $r_i=\sum_cw_{ic}$.
Put
\[
 d_i=\binom{r_i}{2}-\sum_c\binom{w_{ic}}2,
 \qquad H_{ij}=r_ir_j-\sum_cw_{ic}w_{jc}.
\]
Thus $d_i$ counts pairs in row $i$ occupying different columns, and $H_{ij}$ counts pairs from rows $i,j$ occupying different columns.
For $\{u,v\}=[3]\setminus\{c\}$, put
\[
 M_{jk,c}=w_{ju}w_{kv}+w_{jv}w_{ku}.
\]
Then
\begin{equation}\label{eq:path-seven-grid-good-count}
 G=\sum_{i<j}d_id_j
   +\sum_{i=1}^3
       \left(d_iH_{jk}+\sum_c\binom{w_{ic}}2M_{jk,c}\right),
\end{equation}
where $j,k$ are the other two rows in each term of the second sum.

To verify the formula, a good subset has row multiplicities $(2,2,0)$ or $(2,1,1)$.
In the first case, each pair must occupy different columns, giving $d_id_j$.
In the second case, if the doubled-row pair occupies different columns, the other two tokens must occupy different columns.
This gives $d_iH_{jk}$.
If the doubled-row pair occupies one column $c$, the other two tokens must occupy the two different columns outside $c$.
This gives the final minor term.

Suppose first that the row counts are $(3,3,1)$.
Let $A,B$ be the size-three row vectors, and let the last token occupy column $c$.
A size-three row has type $111$, $210$, or $300$, up to column permutation.
Formula~\eqref{eq:path-seven-grid-good-count} becomes
\[
\begin{split}
 G={}&d_Ad_B+d_A(3-B_c)+d_B(3-A_c)\\
 &+\sum_{i\ne c}\binom{A_i}2B_{j(i)}
  +\sum_{i\ne c}\binom{B_i}2A_{j(i)},
\end{split}
\]
where $j(i)$ is the third column different from $i,c$.
The following bounds hold for all relative column placements:
\[
\begin{array}{c|c|c}
 A\text{ type}&B\text{ type}&G\text{ at most}\\\hline
111&111&21\\
111&210&20\\
111&300&12\\
210&210&20\\
210&300&15\\
300&300&18
\end{array}
\]
For $111,210$, the first three terms are $19-3B_c$, and the remaining terms total at most one.
For $111,300$, they are at most nine and three, respectively.
For $210,210$, they are at most sixteen and four.
For $210,300$, they are at most six and nine.
For $300,300$, only the last two terms remain, each at most nine.
The $111,111$ value follows directly.
Thus this row profile has at most twenty-one good subsets.

Now let the row counts be $(3,2,2)$, with row vectors $A,B,C$.
Write $d=d_A$.
Each size-two row is either split between two columns or doubled in one column.

First suppose both $B,C$ are split.
If their missing column is the same column $c$, the formula reduces to
\[
 G=7+4d+2A_c+2\binom{A_c}2\le21.
\]
For $A_c=0,1,2,3$, the largest possible values are $15,21,21,19$, respectively.
If their missing columns differ, let $c$ be their common occupied column.
Their dot product is one, every complementary minor is one, and the formula gives
\[
 G=13+4d-A_c\le24.
\]
Indeed, $d\le2$ gives $G\le21$, while $d=3$ forces $A=(1,1,1)$ and gives $G=24$.

Next suppose $B$ is doubled in column $b$ and $C$ is split.
Using $\sum_i\binom{A_i}2=3-d$ and bounding the complementary minor by $3-A_b$ gives
\[
 G\le15+d(3-2C_b)-3A_b\le21.
\]
For $d\le2$, the last bound is immediate.
For $d=3$, one has $A_b=1$, and the same bound follows.

Finally suppose both $B,C$ are doubled.
If their occupied column is the same, every four-subset is bad, so $G=0$.
Otherwise let their columns be $b,c$, and let $a$ be the remaining column.
The formula gives
\[
 G=4\left(d+\binom{A_a}2+A_a\right)
  =4\left(3-\binom{A_b}2-\binom{A_c}2+A_a\right)\le24.
\]
This completes the upper bound on $G$.

For sharpness, use three matching cells with token counts $3,2,2$.
A four-subset is bad exactly when it uses at most two of those cells.
The count is $\binom54+\binom54+\binom44=11$.
\end{proof}

\begin{theorem}[Missing sectors in a nonplanar seven-pole]
\label{thm:path-seven-eleven-missing-sectors}
Let $X$ be a nonplanar normalized seven-pole of rank $h$.
At least eleven distinct four-port forest totals satisfy
\[
 F_S(X)\le2^h-1.
\]
\end{theorem}

\begin{proof}
Choose a $K_{3,3}$ subdivision and assign all ports to its nine core paths as in Theorem~\ref{thm:path-six-full-four-planar}.
Regard these assignments as seven distinct tokens in the corresponding $3\times3$ grid.
The complement of any core square consists of the paths in one row-column union.
For every four-port subset counted by Lemma~\ref{lem:path-seven-grid-eleven}, choose the corresponding square.
Those four ports lie in one component after deleting its subdivided cycle.
A parity join in that component, together with the cycle, gives a nonforest state in that four-terminal affine class.
There are at least eleven distinct such classes.
Each has $2^h$ states in total, so its forest total is at most $2^h-1$.
\end{proof}

For a seven-pole of girth at least seven, Euler's inequality gives $n\le7$ if the graph is planar.
Since $n=2h+5$, a planar example has rank one and is the seven-cycle.
Thus every such shore of rank at least two is nonplanar and has the eleven missing sectors above.
There are thirty-five four-terminal sectors and seven six-terminal sectors in total.
Consequently, at most thirty-one of these forty-two forest sectors can be full.
This is the count required by the bounded-shore quantum argument at threshold thirty-two.
